\section{Computable relevance zones}
\label{sec:computable-zones}

The soundness theorem gives an explicit ranked zone at each AND node.  Its
four terms use separate clocks for obligation occupancy and defender
completion.

\begin{theorem}[Ranked zone reading]
\label{thm:T4}
\upshape
At an ordinary AND node, the verifier zone is
\[
  Z(N)=Z_{\rm dir}(N)\cup Z_{\rm seed}(N)
       \cup Z_{\rm touch}(N)\cup Z_{\rm virgin}(N).
\]
Any independently nonempty searched set $S(N)\supseteq Z(N)$, together with
\cref{zone:Z4}, is sufficient.  If $Z(N)$ is empty, one arbitrary legal reply
must still be searched.  The uniform admissible deadline $r_N(y)=B(N)$ gives
the radius-$8(B(N)-1)$ seed band in \cref{def:D11}.  Exact cell ranks can only
reduce that band, and exact window exposures can reduce the corresponding
$B$-clock completion guards.  Larger admissible upper bounds remain sound.
Current threat-hitting cells are optional search candidates.  They enter only
the separate kernel of \cref{thm:T6}.
\end{theorem}

\begin{proof}
The four terms and independent fallback are exactly the requirements of
\cref{def:D11}.  Apply \cref{thm:T3}.
\end{proof}

For a short threat-creating continuation, early attacker obligations remain
inside windows already touched by current attacker stones.

\begin{lemma}[Short-range attacker-obligation containment]
\label{lem:L10}
Suppose all attacker placements in a certificate are threat-creating.  If a
placement is the attacker's $k$th future placement from $N$, then for $k\le3$
its window contains at least
\[
  3-(k-1)\ge1
\]
attacker stones already present in $P_N$.  Its obligation role therefore lies
in an attacker-touched window.
\end{lemma}

\begin{proof}
Immediately before the placement, its window has at least three attacker
stones.  At most $k-1$ are earlier future placements, leaving at least
$3-(k-1)$ current stones.
\end{proof}

From the fourth future placement onward, a setup cell can lie in a currently
virgin window and must be supplied by the obligation set.

LOSS leaves need only a small subfamily of their available threat witnesses.

\begin{lemma}[Sparse LOSS witnesses]
\label{lem:L13}
A LOSS witness family may be chosen with $|\mathcal T|\le3$ when $b=1$ and
$|\mathcal T|\le6$ when $b=2$, where $b$ counts placements.
\end{lemma}

\begin{proof}
Every threat empty-set has size one or two.  For $b=1$, if the family contains
$\{a\}$, select it and one set missing $a$.  Otherwise select
$E=\{a,b\}$, one set missing $a$, and one missing $b$.  No point meets all
selected sets.  The three edges of a triangle show sharpness.

For $b=2$, take a maximal pairwise-disjoint subfamily.  It cannot have size
one, since its sole set would meet the full family.  If it has size at least
three, three disjoint sets suffice.  Otherwise take two disjoint sets
$E_1,E_2$.  Their two-point transversals are the at most four cross-pairs
choosing one point from each.  For every cross-pair, select one original set
that it misses.  Together with $E_1,E_2$, at most six sets exclude every
two-point transversal.  The six edges of $K_4$ show that this general
rank-two bound is sharp.  Since a LOSS remainder has $|H|\le b$, a selected
family with transversal number greater than $b$ retains a witness disjoint
from $H$.
\end{proof}

A common radius-three generator covers the full zone only under joint local
budget and witness-location hypotheses.

\begin{theorem}[Static-set coverage for short local budgets]
\label{thm:T5}
\upshape
Suppose $B(N)\le3$, there are at most three remaining attacker placements and
all are threat-creating, every named witness-empty role lies in an
attacker-touched alive window, and $Z_{\rm seed}(N)$ is empty.  Then
\[
  r_3(P_N)\cup
  \{\textup{empties of attacker-touched alive windows at }P_N\},
\]
where
\[
  r_3(P)=\{c\in\Legal(P):\exists s\in\Stones(P),\ d(c,s)\le3\},
\]
covers $Z(N)$.  Without all the joint hypotheses, this static set is a
heuristic candidate generator rather than a certified zone.
\end{theorem}

\begin{proof}
Since $E_N^\D(W)\le B(N)\le3$, the virgin term is empty.  A touched window in
$Z_{\rm touch}$ has at least $6-E_N^\D(W)\ge3$ defender stones, so all its
empties lie within distance $3$ of a defender stone by \cref{lem:L1}.
The previous lemma covers the stated attacker-placement roles, the witness
hypothesis covers the remaining direct roles, and the seed term is empty.
\end{proof}

In a fully forced region, a residual-transversal kernel can replace the
ranked zone while preserving the same path bound.

For the current attacker-threat empty-set family $\mathcal F_N$, write
\[
  \mathcal F_N\setminus d=\{E\in\mathcal F_N:d\notin E\},\qquad
  K_b(N)=\{d\in\Legal(P_N):\tau(\mathcal F_N\setminus d)\le b-1\}.
\]

\begin{theorem}[Extendable-hit kernel]
\label{thm:T6}
\upshape
A kernel-governed region is entered at a certificate node where the real and
ghost positions are equal.  It uses the following rule until a typed terminal
or a handoff with $\mhs>b$.  At every internal AND node in the region, require
the explicit defender $\neg\ownwinnow$ check and
\[
  \mhs(P_N)=\tau(\mathcal F_N)\le b,
\]
and search exactly $K_b(N)$.  No original obligation or core term is required.
If $\mhs<b$, then $K_b(N)=\Legal(P_N)$ and there is no pruning.  If
$\mhs=b$, the kernel weakly refines the current hitting-cell set and can
strictly reduce it; equality can occur.  A node with $\mhs>b$ ends the kernel
region.  It must be a valid LOSS leaf or hand off to an existing nonempty
sound subtree.  The core-free claim does not apply below that handoff without
a fresh equal-position entry.
\end{theorem}

\begin{proof}
If $\tau<b$, deleting the sets hit by any legal $d$ cannot increase $\tau$, so
every legal reply lies in $K_b$.  If $\tau=b$, follow a minimum transversal.
Its first cell is legal by \cref{thm:T1} and lies in $K_b$, so the kernel is
nonempty.

Before the first real reply $d\notin K_b$, the real and certificate plays
agree and $X=Y=\varnothing$.  Abandon the original subtree and use the
residual current threat family.  For $b=2$, the successor has defender budget
$1$ and residual transversal number greater than $1$, giving an adaptive LOSS
refutation.  For $b=1$, a current attacker threat survives and the attacker
has budget $2$, giving a WIN refutation.  No original future obligation is
used.

The defender cannot complete first.  The explicit immediate-win premise is
false.  It bounds every defender-alive window by count $3$ at $b=2$ and count
$4$ at $b=1$.  The remaining two or one defender placements reach at most
$5$.

For the same-$T$ claim, follow a minimum hitting line through the kernel.  At
$b=2$, its first member leaves a family hit by the second, and the
kernel-governed $b=1$ child searches that second member.  At $b=1$, the sole
common hit is searched.  These placements kill every current attacker threat,
and defender stones create no new attacker threat.  The first following
attacker placement cannot complete a window: such a window would have been a
current count-$5$ threat and was killed.  The corresponding original path
therefore reaches the second attacker placement required by the auxiliary
refutation.

Finally, if $\tau>b$ and $d\in K_b$, then $d$ together with a residual
transversal of size at most $b-1$ would meet all of $\mathcal F_N$, a
contradiction.  Thus $K_b$ is empty outside the stated scope and would violate
the independent nonemptiness rule.
\end{proof}

\subsection{Ranked relevance closure}

The mandatory closure is the ranked zone itself.  Search may add local
heuristics without making them verifier premises.

\begin{definition}[Ranked relevance closure]
\label{def:D13}
Define
\[
  R_{\rm cert}(\mathcal C,N)
  =Z_{\rm dir}(N)\cup Z_{\rm seed}(N)
   \cup Z_{\rm touch}(N)\cup Z_{\rm virgin}(N).
\]
For search, the solver may use the finite superset
\[
R_{\rm search}(\mathcal C,N)=R_{\rm cert}(\mathcal C,N)
 \cup\operatorname{hitting}(P_N)\cup\mathcal A(P_N)\cup r_3(P_N),
\]
where $\operatorname{hitting}(P)$ is the union of the current attacker-threat
empty sets, $\mathcal A(P)$ is the set of empties in attacker-touched alive
windows, and $r_3(P)$ is the legal radius-$3$ neighborhood of current stones.
The last three terms are heuristics, not hypotheses of \cref{thm:T3}.
Obligations are finite because the certificate is finite, and touched windows
are finite because the position is finite.  For each finite legal candidate,
the virgin test examines only windows within its bounded exposure radius.
\end{definition}

The verifier must reconstruct every role and clock used by the mandatory
closure.

\begin{theorem}[Ranked-closure coverage]
\label{thm:T7}
\upshape
At every ordinary zone-governed AND node, any independently nonempty
$S(N)\supseteq R_{\rm cert}(\mathcal C,N)$ satisfies \cref{zone:Z2} and
\hyperref[zone:Z5prime]{clause \textup{(Z5$'$)}}.  Dismissal soundness follows from \cref{thm:T3}
when the certificate grammar, the budget, rank, and exposure checks, and
\cref{zone:Z4} also hold.  In constructing $R_{\rm cert}$, the verifier unions
roles over all reachable descendants, includes every designated attacker move
including OR-COMPLETION, treats WIN and LOSS continuation cells as leaf-entry
witness-empty roles, takes the maximum rank over live roles sharing a cell,
and forms no band after the $r=0$ deadline check.  Defender-completion windows
remain on the $E^\D$ clock or the conservative $B$ clock.
\end{theorem}

\begin{proof}
These requirements are exactly the four components and three live verifier
labels of \cref{def:D11}.  Apply \cref{thm:T4,thm:T3}.
\end{proof}

\subsection{Branch-indexed substitution}

A dismissal may select one searched child and verify a smaller envelope over
that child's reachable descendants.  The current transition must be included
in every clock.

\begin{definition}[Transition-inclusive substitution envelope]
\label{def:D17}
At an AND node $N$, annotate each ghost-legal dismissal $d$ with a searched
substitute
\[
  \varphi_N(d)=s\in S(N)
\]
whose exact child is $C_s$.  Declare a nonempty fallback
$F(N)\subseteq S(N)$ independently of these annotations; in particular,
$S(N)$ is not defined as only the replies with no safe substitute.  A valid
selected envelope has the following data and checks.
\begin{enumerate}[label=\textup{(\arabic*)}]\upshape
  \item Its transition budget is
    $\widehat B(N,d,s)=1+B(C_s)$, including the current real/ghost defender
    ply and every reachable LOSS remainder.
  \item Its obligations are the union of roles over all reachable descendants
    of $C_s$, not one leaf or selected continuation.
  \item The real cell $d$ avoids those obligation cells and every
    transition-dangerous completion empty.
  \item For a child obligation $y$ with rank $r_{C_s}(y)$ and
    $y\notin\Legal(P_N)\cup\Stones(P_N)$, direct occupation $d=y$ is forbidden
    and the parent seed test is
    \[
      d(d,y)>8r_{C_s}(y)
      =8\bigl((1+r_{C_s}(y))-1\bigr).
    \]
  \item For each defender-alive touched window $W$ with
    $\cnt_\D(W,P_N)\ge1$, the $\widehat B$ form forbids
    $d\in E(W,P_N)$ when
    $\cnt_\D(W,P_N)+1+B(C_s)\ge6$.  For an all-empty $W$, it forbids a seed
    with $d(d,W)\le8(\widehat B-6)$ when $\widehat B\ge6$.  The exact exposure
    form instead sets
    $\widehat E^\D(W)=1+E_{C_s}^\D(W)$, uses
    $\cnt_\D(W,P_N)+\widehat E^\D(W)\ge6$ for a touched fill, and uses
    $d(d,W)\le8(\widehat E^\D(W)-6)$ with $\widehat E^\D(W)\ge6$ for a virgin
    window.  Both touched and virgin clauses are mandatory for the selected
    clock.
  \item The fallback $F(N)$, and hence $S(N)$, remains nonempty independently
    of whether any dismissal has a safe substitute.
  \item For a ghost-legal A3 dismissal, the ghost plays $\varphi_N(d)$.  Case
    A2 may use any searched filler.  A ghost-illegal A3 dismissal inherits the
    envelope of the earlier ghost-legal seed supporting its real legality and
    may use any searched filler in the current reachable subtree; it does not
    start a child-only clock.
  \item The envelope protects every reachable LOSS witness-empty role through
    leaf entry, while $\widehat B$ or $\widehat E^\D$ counts the leaf's $b$
    placements for defender completion and defender-own-win exclusion.
\end{enumerate}
The default-child form declares $f(N)\in F(N)$ and uses $C_{f(N)}$ for every
ghost-legal dismissal.  It must satisfy the same transition-inclusive tests;
a child obligation union alone is insufficient.
\end{definition}

These envelopes preserve the coupling because their reachable sets remain
nested after the substitution.

\begin{theorem}[Branch-indexed substitution soundness]
\label{thm:T9}
\upshape
The conclusion of \cref{thm:T3} remains valid if an AND node's global
dismissal tests are replaced, dismissal by dismissal, by valid
\cref{def:D17} envelopes.  The same holds for the default-child form.
\end{theorem}

\begin{proof}
The real placement $d$ and ghost substitute $s$ consume the same defender
ply, so clock, mover, and budget remain synchronized.  The avoidance check
supplies occupancy safety.  The seed and completion checks supply legality-
frontier and window-mask safety with the current transition included.
Thereafter all ghost nodes are reachable descendants of $C_s$.  Case A1 may
take any searched edge there; A2 adds no $X$-stone and may take any filler.
At a later ghost-legal A3 step, the selected descendant set nests inside the
earlier one.  A ghost-illegal A3 step traces its legality chain to an earlier
tested seed and inherits that envelope.  Every earlier $X$-stone therefore
continues to avoid current roles, \cref{lem:L9prime} applies with the
transition rank, and \cref{lem:L12} applies with $\widehat B$ or
$\widehat E^\D$.  The LOSS check covers the post-leaf portion.  Attacker
legality, compressed witness transfer, and all terminal cases of
\cref{thm:T3} are unchanged.  The real play either wins earlier or resolves
on the selected reachable path.
\end{proof}

The added transition unit is necessary in both channels.  For the legality
frontier, let $N$ have $b=2$, let legal $d$ be at distance $8$ from a
ghost-illegal future obligation $y$, and let $C_s$ have one defender placement
before a shared legal attacker setup makes designated $y$ legal.  Then
$r_{C_s}(y)=1$.  A child-only radius $8(r-1)=0$ permits $d$; the real move
legalizes $y$ and the remaining defender placement occupies it.  The parent
radius $8r=8$ forbids $d$.  For completion, let a defender-alive window have
two defender stones at $N$, with $d\in W$, $s\notin W$, and $B(C_s)=3$.
The child-only count is $2+3=5$, but real $d$ and three later fills complete
$W$.  The transition-inclusive count is $2+1+3=6$ and forbids the dismissal.

The ordinary completion zone also explains why the internal immediate-win
check is diagnostic rather than an extra hypothesis.

\begin{lemma}[Internal own-win diagnostic]
\label{lem:L14}
Under the completion-zone requirement of \cref{def:D11} and the ban in
\cref{def:D9} on defender-terminal edges, an internal AND node cannot have a
defender-alive count-$5$ window, or a count-$4$ window when $b=2$.
\end{lemma}

\begin{proof}
A count-$5$ touched window has $E_N^\D(W)\ge1$, so its legal last empty lies
in $Z_{\rm touch}$ and must be searched.  Its exact child is
defender-terminal, which is forbidden.  At $b=2$, a count-$4$ touched window
has exposure at least two before the attacker can move, so both empties are
searched.  After either first fill, the $b=1$ child has count $5$.  It cannot
be a LOSS leaf because the mandatory leaf check fails; if it is internal, its
last empty is searched and gives the same forbidden terminal edge.
\end{proof}

The diagnostic remains mandatory at LOSS leaves and an explicit premise of
the kernel theorem, whose region has no completion guard.

\subsection{Finite certificate DAGs}

Shared certificate subtrees are sound when all shared labels and clocks are
exactly consistent.

\begin{definition}[Finite acyclic certificate DAG]
\label{def:D18}
A certificate may be a finite rooted acyclic graph if every shared node has
one exact \cref{def:D9} label: the same position, mover and budget, designated
OR action or AND searched-successor map, leaf witness data, and one consistent
path clock.  Obligations at a DAG node are unions over all reachable
descendants.  Local budgets, ranks, and exposures satisfy their inequalities
on every outgoing edge.  Coupling histories $X$, $Y$, and $\widehat X$ remain
path-local.
\end{definition}

Finite unfolding reduces the DAG form to the tree theorem.

\begin{theorem}[DAG unfolding]
\label{thm:T10}
\upshape
Every \cref{def:D18} DAG satisfying the ordinary \cref{def:D11} rules or the
substitution rules of \cref{def:D17} has the corresponding conclusion of
\cref{thm:T3}.
\end{theorem}

\begin{proof}
A finite acyclic graph has finitely many root-to-node paths.  Unfold one copy
of each shared node along each such path.  Exact labels and consistent clocks
give a finite \cref{def:D9} tree with the same legal transitions, terminal
data, and local-clock inequalities.  Reachability remains nested along every
edge, so descendant obligation unions and protection monotonicity are
preserved.  Apply \cref{thm:T3} or \cref{thm:T9}.  The unfolded coupling
histories are exactly the required path-local histories.
\end{proof}

The solver may search with $R_{\rm search}$, a smaller heuristic, or a
substitution-aware generator.  Verification must establish ordinary
$R_{\rm cert}$ coverage and attacker legality, or validate every substitution
annotation, while retaining an independently nonempty fallback.  A failure
yields \textsc{unknown}.  Search trimming affects completeness only.

The predecessor heuristic measured a mean of $147$ cells against $302$ legal
cells on random midgame positions.  These figures describe the conservative
search generator, not the ranked certificate-relative verifier zone.
